%\section{Synopsis}
%Thermodynamic modeling of geological fluids is essential for the design of geothermal energy and \ce{CO2} sequestration technologies

\section{Introduction}

In the global warming context caused partly by the rising level of anthropogenic \ce{CO2} in the atmosphere, alternatives to non-renewable fossil energy, such as oil, coal, and natural gas, are seriously pursued. Geothermal energy stands out as a source of low-carbon sustainable energy. Compared to other renewable energies, e.g. biomass, wind, solar, and tidal, geothermal presents lower operation costs, higher stability, and higher utilization rate \cite{Wu2020}. The operation consists of extracting a hot fluid from a geothermal subsurface formation and converting its thermal energy to electric power. Geothermal formations are either saturated with fluids or hot dry rocks (HDR), which are fractured to increase their permeability and improve heat extraction efficiency. Water is usually the main component in natural geothermal formations. In some superhot geothermal formations, the concentration of \ce{CO2} may be very high, reaching up to 7\% weight fraction \cite{Armannsson2017}. Geothermal formations may have a wide range of temperatures, from 100 ºC to 400 ºC or higher depending on the geothermal gradient \cite{Dickson2003}. 

Recently, \ce{CO2} has been suggested as a heat carrier fluid, which is a technique from Enhanced Geothermal Systems (EGS).
With \ce{CO2} as the working fluid, the heat extraction efficiency in fractured HDR formations may increase.  A large amount of the working fluid remains underground during the operation, which in the case of \ce{CO2}-EGS represents a potential to couple carbon sequestration (CCS) and geothermal energy extraction. Unlike other CCS techniques, the \ce{CO2} sequestration through EGS is associated with power production, which compensates for the cost of geological carbon storage. \ce{CO2} is equivalent to water as a heat carrier fluid and may have promising thermodynamic and transport properties (e.g. lower viscosity, higher compressibility, and expansibility) that facilitate the operation. For instance, the significant difference between the density in the injection (cold) and production (hot) wells provides a large buoyancy difference that reduces power requirements to circulate the fluid \cite{Wu2020}. From the hydraulic fracturing by water, there may be a large amount of water in the geothermal formation, either due to the fracturing itself or from the formation brine \cite{Wu2020}. Consequently, the resulting working fluid is a mixture of \ce{CO2}-water and \ce{CO2}-brine. The knowledge of thermophysical properties, including phase equilibria and density, of the \ce{CO2}-water and \ce{CO2}-brine mixture at geothermal conditions is essential to simulate and optimize EGS processes. At lower temperatures, other applications in which knowledge of phase behavior between \ce{CO2} and water is important includes CCS in deep saline aquifers, enhanced oil recovery (EOR), \ce{CO2} in the food industry (e.g. carbonated hydroalcoholic drinks) and \ce{CO2}-supercritical extraction in the chemical industry.

Equations of state (EoS) can be used to model the behavior of \ce{CO2} and water mixtures from a thermodynamic perspective. \citeauthor{Zhao2016} \cite{Zhao2016} have described \ce{CO2}-\ce{H2O} phase equilibria with high accuracy for a wide range of pressure (0.1 - 200 MPa) and temperature (273 - 623 K) coupling a cubic EoS (Peng-Robinson-Stryjek-Vera) with a mixing rule from an excess Gibbs free energy ($G^{E}$) model (Wong-Sandler with NRTL). To obtain agreement with experiments, four adjustable temperature-dependent parameters are required, and the binary interaction parameter ($k_{ij}$) is made composition-dependent. \citeauthor{Pappa2009} \cite{Pappa2009} also have applied the cubic Peng-Robinson EoS to describe the phase equilibria. A binary interaction parameter is used not only in the cross-energy parameter ($a_{ij}$) but also in the cross co-volume parameter ($b_{ij}$), which are both made temperature-dependent but with two different regimes: at temperatures lower than 373 K a set of equations is used to describe the interaction parameters, and at higher temperatures a different set of equations is applied \cite{Pappa2009}.

Water molecules self-associate with other water molecules through hydrogen bonding. To have a model physically consistent, one must explicitly account for association in aqueous mixtures. The system free energy due to highly directional interactions can be obtained from Wertheim's theory \cite{Wertheim1984}. The Cubic-Plus-Association (CPA) EoS \cite{Kontogeorgis1996} combines dispersion interaction described from cubic EoS to association described from Wertheim's theory. If there is no associating species in the mixture, the cubic EoS is recovered. CPA EoS has been applied to describe \ce{CO2}-\ce{H2O} phase equilibria \cite{Tsivintzelis2011, Li2009, Pappa2009, Monteiro2020, Perfetti2008}. Water molecules are commonly modeled with a 4C scheme (2 electron-donor and 2 proton-donor sites). Different association schemes for \ce{CO2}, including no association, have been proposed \cite{Tsivintzelis2011}. Higher-accuracy predictions are obtained with the \ce{CO2}-\ce{H2O} cross-association through solvation and without \ce{CO2} self-association \cite{Tsivintzelis2011, Li2009, Xiong2023, Bian2019}. Although good agreement is obtained, previous CPA models have limited the maximum temperature to about 530 K \cite{Li2009, Tsivintzelis2011}. \textcolor{red}{\citeauthor{Perfetti2008} \cite{Perfetti2008} have included dipolar interactions in the CPA EoS to describe the \ce{CO2}-\ce{H2O} phase equilibria; however, the critical pressures of the mixture at temperatures higher than 500 K are overestimated.} The Statistical Associating Fluid Theory (SAFT) EoS \cite{Chapman1989} can also be used to model the phase equilibria from a molecular perspective. \citeauthor{Diamontonis2012} \cite{Diamontonis2012} have modeled \ce{CO2} as associating component with the Perturbed Chain-SAFT (PC-SAFT) and as a quadrupolar component with the truncated PC-Polar SAFT (tPC-PSAFT). By comparing the predictions, the \ce{CO2}-\ce{H2O} cross-association scheme outperforms the others. The model accurately predicts the vapor-liquid (VLE), liquid-liquid (LLE) equilibria and saturated liquid density, but the maximum temperature is limited to about 530 K \cite{Diamontonis2012}. \textcolor{red}{SAFT-$\gamma$ Mie has also been applied to model \ce{CO2}-water phase equilibria and Henry's constant \cite{Papaioannou2016}. Accurate results with fewer mixture parameters have been obtained, but a limited range of conditions was examined.}

When there are electrolytes in the mixture, the thermodynamic model must account for electrostatic interactions. Simplistic models could represent reasonably the phase equilibria via numerical correlations \cite{Sun2021}
or classical EoS \cite{Bian2019, Yang2024} by adding several unphysical adjustable parameters. The application of such models is limited exclusively in the range they are derived. Electrolytes decrease the \ce{CO2} solubility in the aqueous mixture significantly; a phenomenon known as the ``salting-out'' effect. \citeauthor{Li2020} \cite{Li2020} introduce the Debye-H\"uckel activity coefficient to correct the chemical potential of nonelectrolyte components in the aqueous phase. The $\gamma-\phi$ (activity coefficient - fugacity coefficient) has also been applied  by \citeauthor{Springer2012} \cite{Springer2012}, in modifying the excess Gibbs energy model with  Pitzer–Debye–H\"uckel equation to account for long-range electrostatic interactions. Several fitting parameters are required in the $\gamma-\phi$ approach, and the model may fail at higher pressure because $G^E$ is pressure-independent.

The electrolyte Cubic-Plus-Association EoS (eCPA) \cite{Maribo2015} is a relatively simple and physically consistent $\phi-\phi$ model that has been used to describe electrolyte solutions. In this approach, a single equation of state describes simultaneously both the aqueous and \ce{CO2}-rich phases. The residual Helmholtz free energy contribution related to ion-ion electrostatic interactions is added to the CPA. It is usually represented by the Debye-Hückel (DH) or mean spherical approximation (MSA) equations. The DH term originates from solving the linearized Poisson-Boltzmann (PB) equation  \cite{Michelsen2004}.  The linear PB assumption holds for dilute solutions and deviations appear at higher ionic strength \cite{Debye1923}. 
%The DH equation was developed assuming a constant relative permittivity. Nevertheless, this constraint should be relaxed when using this model.
By solving the PB equation, a self-potential known as the Born term arises \cite{Debye1923}. The Born term is related to the free energy variation from the ion charging, accounting for its transfer from vacuum to a medium with a dielectric constant \cite{Michelsen2004}. The Born term accounts for the interaction between ions and the solvent and should be included in the eCPA to keep it physically consistent. Both DH and Born terms require a model for the mixture dielectric constant. When a primitive model is selected, the relative permittivity is estimated from the association term of the eCPA EoS, which provides information on dipole-dipole correlations \cite{Maribo2013, Maribo2013a}. 
%The eCPA parametrization is a key aspect of its success and it can be made both ion \cite{Sun2019} or salt \cite{Olsen2021} specific. The latter may be more appropriate for mixtures with a single salt, because a reduced number of parameters is required. Nevertheless, for mixtures with multiple ionic species, the system may become over-specified and an ion-specific approach is more adequate \cite{Maribo2015}. The target properties in the parametrization are also important. By including brine density data in the eCPA parametrization, the density can be predicted without volume translation, but on the other hand, the dielectric saturation (decrease in the dielectric constant with salinity) is not captured \cite{Olsen2021}. 

The \ce{CO2}-brine phase equilibria have been described with the eCPA EoS \cite{Olsen2021, Chabab2019, Sun2019, Xiong2023, Maribo2015}. In most cases, \ce{CO2}-\ce{H2O} cross-association is not accounted for, which leads to larger deviations and inconsistent \ce{CO2} solubility trend with temperature \cite{Maribo2015, Olsen2021}. To reproduce the experimental data, \citeauthor{Chabab2019} \cite{Chabab2019} applied an eCPA with a non-primitive model for the dielectric constant and with binary interaction parameters dependent on temperature, salt concentration, and phase state. By including cross-association, \citeauthor{Xiong2023} \cite{Xiong2023} have improved the \textcolor{red}{LLE} and \textcolor{red}{VLE}  predictions. They have assumed, however, that the association strength is the same as the molar volume to simplify the relative permittivity model and have limited the temperature range from 298 to 473 K \cite{Xiong2023}. \textcolor{red}{\citeauthor{Zoghi2012} \cite{Zoghi2012} have successfully modeled the solubility of \ce{CO2} in aqueous N-methyldiethanolamine (MDEA) in a wide range of pressures, temperatures, acid gas loadings, and MDEA concentrations using a modified Peng-Robinson eCPA EoS that accounts for short- and long-range ionic interactions with the MSA theory.}

Phase equilibria of \ce{CO2}-water and \ce{CO2}-brine have been investigated by Molecular Dynamics (MD) \cite{Liu2013, Blazquez2024} and Monte Carlo (MC) \cite{Vorholz2000, Vorholz2004, Liu2011, Jiang2017a, Orozco2014, Vlcek2011} simulations. From the direct coexistence MD method, which provides a description of the phase interface, the interfacial tension is obtained in addition to each phase composition and density. The large free energy of the interface may represent a diffusion barrier and the final composition may not be at equilibrium when a large enough system size and simulation length are not selected. From histogram-reweighting grand canonical Monte Carlo, one can obtain several points of the phase diagram from the fluctuations of a single simulation \cite{Panagiotopoulos2000}. In the Gibbs ensemble Monte Carlo (GEMC) simulations, two independent boxes are simulated simultaneously. For a given pressure and temperature, the composition and volume of each box may change until equilibrium is reached \cite{Panagiotopoulos1992}. 

The accuracy of the molecular simulations rely on the force field (FF). \citeauthor{Liu2011} \cite{Liu2011} evaluate various combinations for  water (SPC, TIP4P, TIP4P2005, and exponential-6) and \ce{CO2} (EPM2, TraPPE, and exponential-6) models. Using Lorentz-Berthelot (LB) combining rules for the unlike parameters, the models do not reproduce adequately the experimental data in the entire temperature range (323 < $T$ < 623 K). The prediction, especially in the aqueous phase, can be improved when optimized unlike parameters are selected \cite{Orozco2014, Vlcek2011}. The simultaneous accurately prediction of both phases, however, requires further considerations \cite{Orozco2014}. The SPCE water model is widely used for the investigation of many properties in aqueous mixtures. Due to the polarization correction, SPCE has a higher dipole moment, which deteriorates the vapor pressure predictions \cite{Vorholz2004}. As a consequence, the solubility of water in nonpolar environments is usually underestimated by SPCE \cite{Vlcek2011}. Using optimized unlike parameters determined at temperatures between 298 and 348 K, \citeauthor{Vlcek2011} \cite{Vlcek2011} have shown that GEMC simulations can predict the equilibrium composition in both phases if self-polarization correction is included for the chemical potential of SPCE water in the \ce{CO2}-rich phase. Even using polarizable force fields for water, accurate predictions in both phases are only achieved if a hydrogen bonding potential is used explicitly to account for charge transfer effects along with optimized cross-interaction parameters \cite{Jiang2017a}. The reduction in the \ce{CO2} solubility in the aqueous phase with salts can also be reproduced through molecular simulations \cite{Blazquez2024, Vorholz2004, Liu2013}. \citeauthor{Blazquez2024} \cite{Blazquez2024} have shown that nominal integer charge for the electrolytes may overestimate the salting-out effect, and scaled-charges can be an alternative. By scaling down ionic charge, one could represent indirectly charge transfer and polarization effects, improving thermophysical property predictions \cite{Zeron2019, Coelho2024}. In most cases, molecular simulations of the \ce{CO2}-\ce{H2O} and \ce{CO2}-brine phase equilibria have been focused on low to moderate temperatures, with few studies \cite{Liu2011} at temperatures higher than 523 K. 

In this work, we evaluate the phase equilibria of \ce{CO2}-water and \ce{CO2}-brine mixtures focusing on geothermal conditions (high temperature) using a multiscale approach. The conditions for \ce{CO2} sequestration in saline aquifers are also included in our investigation. We derive a CPA and an eCPA EoS accounting for \ce{CO2}-\ce{H2O} cross-association to model simultaneously the phase equilibria and density. We compare the results with GEMC simulations using various combining rules. The association in the system is analyzed by both equation of state and molecular simulations.
